\documentclass{article}
\usepackage{pathfinder-readable}

\title{When a Donor Wrist Cannot Locate a Held Object}

\begin{document}
\maketitle

\begin{abstract}
This account reads Aero Hand Open and DynaMAC together at an occluded handover. The thread asked
whether the donor arm's wrist pose can stand in for the pose of an object that its hand rotates
within the grasp. The papers and the peers' checks establish a conditional information limit: a
fixed wrist does not fix the orientation of such an object, and identical available observations
cannot distinguish two orientations that call for different receiving approaches. They do not
establish that this ambiguity occurs often on the physical hand or that DynaMAC fails there. The
verifier therefore left the thread at DRAFT, with a controlled pose measurement and handover
comparison as the next step.
\end{abstract}

\section{Introduction}

Aero Hand Open, paper Q, presents a tendon-driven hand with sixteen joints driven by seven motors
\cite{q}. Its only proprioception comes from the motor encoders. The paper supplies a simulation
model and a learned policy that rotates a cube within the grasp. It reports rotation in simulation
and on the physical hand, but does not log the cube's physical pose during that run.

DynaMAC, paper P, learns manipulation from demonstrations using several reference frames, including
object frames and the other arm's end-effector frame \cite{p}. It selects useful streams as a task
unfolds and uses a moving arm as a reference for bimanual work. For a handover, P says that the held
object's pose and the carrying arm's pose provide redundant information to the receiver. Since the
object can be hidden from view, P proposes the other arm's forward-kinematic pose as a substitute.

The thread pursued that last claim in the particular case where Q's hand carries and rotates the
object. The question was whether the carrying wrist still gives the receiver enough information
about the hidden object. This is a narrower question than installing Q's hand in P's benchmark: the
papers use different simulators and action spaces, so that change would require substantial
integration.

\section{What was done}

The peers first considered a broad swap of Q's hand into P's bimanual tasks. They then examined P's
kinematic-link detector. An early suggestion was that motion within a compliant grasp would make the
detector retain an object stream that should be removed. On checking P's definition, both peers
withdrew that inference. Its geometric mean standard deviation measures the spread of a learned
policy stream across demonstrations at a given phase. It is not a measurement of object motion
relative to the wrist during one execution. A repeatable in-hand movement could therefore leave that
stream precise.

The peers then focused on P's handover argument. They compared P's claim about the donor arm's pose
with Q's demonstrated in-hand cube rotation. They formulated a matched-observation test, checked the
limits of the proposed inference, and proposed controls that would expose them. A prior-work check
found in-hand tracking under occlusion by Liang and colleagues \cite{liang2020} and pose-aware
bimanual handover work by Vezzani and colleagues \cite{vezzani2017}. The ledger does not contain an
exhaustive novelty search.

\section{Results}

Let \(T_o\) be the object's pose, \(T_w\) the carrying wrist's pose, and \(G\) the object's pose relative to that wrist. Then
\[
T_o = T_w G.
\]
The wrist is a sufficient stand-in for the object only to the extent that \(G\) stays fixed, or
changes too little to matter to the receiving task. P's handover discussion relies on this practical
redundancy \cite{p}. Q's rotation task supplies a concrete way for \(G\) to change while the wrist
remains still \cite{q}. The peers identified this as a limit on the wrist substitution, not as an
observed failure of P's system.

The stronger statement is conditional. Suppose two trials give the receiver the same wrist pose and
the same histories of motor readings and actions, but the hidden object has different orientations.
Suppose also that those orientations require different receiving approaches. A receiver given only
those histories has the same input in both trials, so it cannot reliably choose the
orientation-specific approach in both. This is the peers' matched-state argument. Reseating a marked
object under occlusion could create such a test; its natural frequency is unknown.

Q does not log physical cube pose or measure the spread of \(G\) at matched motor readings \cite{q}.
Neither paper reports DynaMAC receiving an object from Q's hand. The result is a possible failure
mode, not a measured handover loss.

\section{Objections and unresolved points}

P's stream spread across demonstrations may or may not alter its link decision under varied grasps
\cite{p}. Detector decisions need comparison with controlled contact states or motor perturbations.
No detector error rate is available.

P omits gripper actions from its pose-level exposition for brevity; that omission does not show that
its implementation cannot command a hand \cite{p}. Q's seven encoders do not, by themselves, prove
that object pose is ambiguous after a known action history. Moreover, if one grasp safely receives
every tested orientation, the proposed orientation distinction has no practical effect. Under
genuinely identical observations and complete occlusion, an uncertainty-aware policy can represent
doubt, seek a new view, or use a robust grasp, but it cannot recover the hidden orientation without
new information.

Object-pose tracking and pose-aware handover precede this thread \cite{liang2020,vezzani2017}. A
controlled test of P's wrist substitution with Q's hand might be distinctive; that remains
unestablished.

\section{Why the thread stopped}

The verifier judged the pose relation and its qualification to be a specific result of reading these
two papers together. It kept the status at DRAFT because the result is a conditional limit, not a
measured DynaMAC failure. The smallest restart is to hold the donor wrist fixed, use a marked
object, and measure \(G\) externally while varying its in-hand orientation or contact state. The
first check is whether materially different object poses remain within matched wrist, encoder, and
action-history conditions. If they do, an orientation-sensitive receiving task can compare
wrist-only inference with visible-object and last-seen-object controls, while checking whether a
single grasp succeeds for both orientations.

\bibliographystyle{plainurl}
\bibliography{references}
\end{document}
